%======================================================================
% Appendix: Backprop-to-EP Transfer Guarantees for Phasor Networks
%======================================================================
% This appendix provides a formal analysis of the equivalence between
% feedforward (backprop) inference and free-phase EP settling in
% phasor networks with soft projection. The results bound the
% difference in hidden states and classification accuracy.
%
% Compile: pdflatex bp_ep_transfer_appendix.tex
%======================================================================

\documentclass[11pt]{article}
\usepackage{amsmath,amssymb,amsthm}
\usepackage{bm}
\usepackage{booktabs}
\usepackage[margin=1in]{geometry}
\usepackage{hyperref}
\usepackage{natbib}
\usepackage{tikz}
\usetikzlibrary{arrows.meta,positioning,calc}

% Macros
\newcommand{\R}{\mathbb{R}}
\newcommand{\C}{\mathbb{C}}
\newcommand{\T}{\mathbb{T}}
\newcommand{\im}{\mathrm{i}}
\newcommand{\dd}{\mathrm{d}}
\newcommand{\dt}{\Delta t}
\newcommand{\unitproject}{\mathrm{unit\_project}}
\DeclareMathOperator{\Real}{Re}
\DeclareMathOperator{\Imag}{Im}
\DeclareMathOperator{\Tr}{Tr}
\DeclareMathOperator{\diag}{diag}
\DeclareMathOperator{\spanop}{span}
\newcommand{\norm}[1]{\left\|#1\right\|}
\newcommand{\abs}[1]{\left|#1\right|}
\newcommand{\ip}[2]{\left\langle #1, #2 \right\rangle}
\newcommand{\softproj}[1]{u_\varepsilon(#1)}

\title{Appendix: Backprop-to-Equilibrium-Propagation Transfer Guarantees\\for Phasor Networks}
\author{}
\date{}

\begin{document}
\maketitle

\begin{abstract}
We analyze the transfer of backpropagation-trained phasor networks to
Equilibrium Propagation (EP) free-phase settling. For a depth-$L$ chain
with soft projection $u_\varepsilon(g) = g/\sqrt{|g|^2+\varepsilon^2}$,
symmetric feedback, and weight scaling $\rho = \max_\ell \|W_\ell\|_2$,
we prove that the EP fixed point $z^*$ and feedforward fixed point
$z^{\text{ff}}$ satisfy $\|z^*_\ell - z^{\text{ff}}_\ell\| \le C_\ell \rho$
for explicit per-layer constants $C_\ell$. This implies a classification
accuracy drop bounded by $O(\rho/\gamma)$ where $\gamma$ is the
logit margin. The result holds for the co-rotating frame
($\texttt{carriers}=\texttt{nothing}$) with $K_\text{mode}=0$.
\end{abstract}

%======================================================================
\section{Introduction and Setup}
%======================================================================

Phasor networks \citep{olin2021deep,olin2023hyperdimensional} are
trained either by backpropagation (feedforward pass) or by
Equilibrium Propagation (iterative settling to a fixed point). A
practical question: \emph{does a backprop-trained network retain its
accuracy when run as an EP network?} Empirically on FashionMNIST
(784$\to$256$\to$64), the drop is $<1\%$ absolute. This appendix
provides a formal bound.

\subsection{Network Architecture}
Consider a depth-$L$ chain of \texttt{PhasorDense} layers:
\begin{align}
  z_0 &= x \in \T^{d_0} \quad\text{(input, fixed)}\\
  z_\ell &= u_\varepsilon\big(W_\ell z_{\ell-1} + b_\ell\big)
  \quad\text{for } \ell=1,\dots,L \quad\text{(feedforward)}
\end{align}
where $W_\ell \in \R^{d_\ell \times d_{\ell-1}}$ are real weights,
$b_\ell \in \C^{d_\ell}$ are complex biases, and
$u_\varepsilon(g) = g/\sqrt{|g|^2+\varepsilon^2}$ is the soft projection
(applied entrywise). The output layer $L$ has no feedback in the
feedforward pass.

\subsection{EP Free-Phase Dynamics}
In EP with $K_\text{mode}=0$, $\beta=0$, symmetric feedback
$B_\ell = W_{\ell+1}^\top$, and co-rotating frame, the settle operator
is:
\begin{align}
  F_\varepsilon(z)_\ell &= u_\varepsilon\big(W_\ell z_{\ell-1} +
  W_{\ell+1}^\top z_{\ell+1} + b_\ell\big), \quad \ell=1,\dots,L-1 \\
  F_\varepsilon(z)_L &= u_\varepsilon\big(W_L z_{L-1} + b_L\big)
\end{align}
with $z_0 = x$ clamped. The free fixed point $z^*$ satisfies
$z^* = F_\varepsilon(z^*)$.

\subsection{Assumptions}
\begin{enumerate}
  \item[(A1)] \textbf{Weight scaling:} $\rho := \max_{\ell=1}^L \|W_\ell\|_2 \le \rho_\text{max}$.
    In experiments, weights are scaled by $0.4$, giving $\rho \approx 0.4$.
  \item[(A2)] \textbf{Bias bound:} $\max_\ell \|b_\ell\| \le B$.
  \item[(A3)] \textbf{Soft projection parameter:} $\varepsilon > 0$ fixed
    (experiments use $\varepsilon = 0.1$).
  \item[(A4)] \textbf{Co-rotating frame:} $\texttt{carriers}=\texttt{nothing}$,
    so lab-frame rotation is exact and does not affect the co-rotating
    dynamics.
  \item[(A5)] \textbf{Contraction condition:} $\rho < \frac{1}{2L}$.
    This ensures $F_\varepsilon$ is a contraction (Lemma 2).
\end{enumerate}

%======================================================================
\section{Soft Projection Properties}
%======================================================================

\begin{lemma}[Soft projection Lipschitz bounds]
\label{lem:soft-lipschitz}
For $u_\varepsilon(g) = g/\sqrt{|g|^2+\varepsilon^2}$ applied entrywise
to vectors in $\C^d$:
\begin{enumerate}
  \item[(i)] $\|u_\varepsilon(a) - u_\varepsilon(b)\|_2 \le \|a - b\|_2$
    (1-Lipschitz in Euclidean norm).
  \item[(ii)] Angular distance on the torus:
    $d_\T(u_\varepsilon(a), u_\varepsilon(b)) \le d_\C(a,b)$.
  \item[(iii)] Jacobian operator norm: $\|Du_\varepsilon(g)\|_\text{op}
    \le \frac{1}{\varepsilon}\big(1 + \frac{|g|}{\sqrt{|g|^2+\varepsilon^2}}
    \big) \le \frac{2}{\varepsilon}$.
\end{enumerate}
\end{lemma}

\begin{proof}
(i) The map $g \mapsto g/\sqrt{|g|^2+\varepsilon^2}$ is the gradient of
the convex function $\phi(g) = \sqrt{|g|^2+\varepsilon^2}$, hence
1-Lipschitz. Entrywise application preserves this.
(ii) Angular distance is bounded by Euclidean distance in the ambient
space. (iii) Direct computation of the derivative.
\end{proof}

\begin{remark}[Hard projection limit]
The hard projection $\unitproject(g) = g/|g|$ is the pointwise limit
$\varepsilon \to 0$ of $u_\varepsilon$. However, $\unitproject$ is
discontinuous at $g=0$ and not Lipschitz near zero. Our bounds hold
uniformly for $\varepsilon > 0$. In practice, weight scaling $\rho
\approx 0.4$ keeps drives away from zero, and the hard-projection
fixed point is well-approximated by the soft fixed point with small
$\varepsilon$ (empirically $\varepsilon=0.1$). A formal $\varepsilon
\to 0$ limit requires excluding neighborhoods of zero drives.
\end{remark}

%======================================================================
\section{Contraction of the EP Operator}
%======================================================================

\begin{lemma}[Contraction of $F_\varepsilon$]
\label{lem:contraction}
Under assumptions (A1)--(A5), the EP settle operator $F_\varepsilon$
is a contraction on $(\T^{d_1} \times \dots \times \T^{d_L})$ with
respect to the Euclidean norm on the ambient space $\C^{d_1} \times
\dots \times \C^{d_L}$. The contraction factor is
$\kappa = 2\rho L < 1$.
\end{lemma}

\begin{proof}
For any $z, z'$, by Lemma \ref{lem:soft-lipschitz}(i):
\begin{align}
  \|F_\varepsilon(z)_\ell - F_\varepsilon(z')_\ell\|
  &\le \|W_\ell(z_{\ell-1}-z'_{\ell-1})
    + W_{\ell+1}^\top(z_{\ell+1}-z'_{\ell+1})\| \\
  &\le \rho \|z_{\ell-1}-z'_{\ell-1}\| + \rho \|z_{\ell+1}-z'_{\ell+1}\|
\end{align}
for $\ell=1,\dots,L-1$, and for $\ell=L$:
$\|F_\varepsilon(z)_L - F_\varepsilon(z')_L\| \le \rho \|z_{L-1}-z'_{L-1}\|$.
Summing over $\ell$:
\begin{align}
  \|F_\varepsilon(z) - F_\varepsilon(z')\|
  &\le \rho \sum_{\ell=1}^{L-1} (\|z_{\ell-1}-z'_{\ell-1}\|
    + \|z_{\ell+1}-z'_{\ell+1}\|) + \rho \|z_{L-1}-z'_{L-1}\| \\
  &\le 2\rho L \max_\ell \|z_\ell - z'_\ell\| \le 2\rho L \|z - z'\|
\end{align}
where the last inequality uses $\|z\|^2 = \sum_\ell \|z_\ell\|^2$.
The factor $2\rho L$ follows from each $z_\ell$ appearing in at most
two terms (as $z_{\ell-1}$ and $z_{\ell+1}$), except boundaries.
\end{proof}

\begin{corollary}[Existence, uniqueness, and convergence]
Under (A5), $F_\varepsilon$ has a unique fixed point $z^*$, and the
iterates $z^{k+1} = F_\varepsilon(z^k)$ converge exponentially:
$\|z^k - z^*\| \le \kappa^k \|z^0 - z^*\|$.
\end{corollary}

%======================================================================
\section{Feedforward Fixed Point}
%======================================================================

The feedforward operator $G_\varepsilon$ (no feedback) is:
\begin{align}
  G_\varepsilon(z)_\ell = u_\varepsilon(W_\ell z_{\ell-1} + b_\ell),
  \quad \ell=1,\dots,L
\end{align}
with $z_0 = x$. Since $G_\varepsilon$ is lower-triangular (each layer
depends only on previous), it has a unique fixed point $z^{\text{ff}}$
computed by a single forward pass:
$z^{\text{ff}}_\ell = u_\varepsilon(W_\ell z^{\text{ff}}_{\ell-1} + b_\ell)$.
This is exactly the backpropagation feedforward pass.

%======================================================================
\section{Perturbation Bound: EP vs Feedforward Fixed Points}
%======================================================================

\begin{theorem}[Per-layer perturbation bound]
\label{thm:perturbation}
Let $z^*$ be the EP fixed point and $z^{\text{ff}}$ the feedforward
fixed point. Define the perturbation $\delta_\ell = z^*_\ell -
z^{\text{ff}}_\ell$. Under assumptions (A1)--(A5), for each layer
$\ell = 1,\dots,L$:
\begin{equation}
  \|\delta_\ell\| \le C_\ell \rho + D_\ell B
\end{equation}
where the constants $C_\ell, D_\ell$ are given by the following
recurrence (with $C_0 = D_0 = 0$):
\begin{align}
  C_\ell &= \frac{\rho}{1-\kappa}\Big(
    C_{\ell-1} + C_{\ell+1} + \|z^{\text{ff}}_{\ell+1}\| \Big) \\
  D_\ell &= \frac{\rho}{1-\kappa}\Big(
    D_{\ell-1} + D_{\ell+1} \Big) + \frac{1}{1-\kappa}
\end{align}
for $\ell=1,\dots,L-1$, and for the output layer $\ell=L$:
\begin{align}
  C_L &= \frac{\rho}{1-\kappa} C_{L-1} \\
  D_L &= \frac{\rho}{1-\kappa} D_{L-1} + \frac{1}{1-\kappa}
\end{align}
with the convention $C_{L+1} = D_{L+1} = \|z^{\text{ff}}_{L+1}\| = 0$.
\end{theorem}

\begin{proof}
The difference between the two fixed-point equations is:
\begin{align}
  \delta_\ell &= u_\varepsilon(W_\ell z^*_{\ell-1} + W_{\ell+1}^\top z^*_{\ell+1} + b_\ell)
    - u_\varepsilon(W_\ell z^{\text{ff}}_{\ell-1} + b_\ell) \\
  &= u_\varepsilon(W_\ell z^{\text{ff}}_{\ell-1} + W_{\ell+1}^\top z^{\text{ff}}_{\ell+1} + b_\ell
    + W_\ell \delta_{\ell-1} + W_{\ell+1}^\top \delta_{\ell+1})
    - u_\varepsilon(W_\ell z^{\text{ff}}_{\ell-1} + b_\ell)
\end{align}
Apply Lemma \ref{lem:soft-lipschitz}(i):
\begin{align}
  \|\delta_\ell\|
  &\le \|W_\ell \delta_{\ell-1} + W_{\ell+1}^\top \delta_{\ell+1}
    + W_{\ell+1}^\top z^{\text{ff}}_{\ell+1}\| \\
  &\le \rho \|\delta_{\ell-1}\| + \rho \|\delta_{\ell+1}\| + \rho \|z^{\text{ff}}_{\ell+1}\|
\end{align}
for $\ell < L$. For $\ell = L$, the feedback term is absent:
$\|\delta_L\| \le \rho \|\delta_{L-1}\|$.

This gives a linear system $\|\delta\| \le \rho A \|\delta\| + \rho v$
where $A$ is the tridiagonal matrix of the chain (with $A_{L,L}=0$)
and $v_\ell = \|z^{\text{ff}}_{\ell+1}\|$. Since $\rho < 1/(2L)$,
$\rho A$ has spectral radius $\le 2\rho L = \kappa < 1$. Thus
$(I - \rho A)$ is invertible with $\|(I - \rho A)^{-1}\| \le 1/(1-\kappa)$.
Solving:
$\|\delta\| \le (I - \rho A)^{-1} \rho v$.

The recurrence for $C_\ell, D_\ell$ follows by separating the terms
proportional to $\rho \|z^{\text{ff}}_{\ell+1}\|$ (feedback-induced)
from the bias terms. The bias contribution $b_\ell$ adds a constant
$1$ in the recurrence because $\|u_\varepsilon(g+b) - u_\varepsilon(g)\|
\le \|b\|$ by Lemma \ref{lem:soft-lipschitz}(i).
\end{proof}

\begin{corollary}[Simplified global bound]
Let $Z = \max_\ell \|z^{\text{ff}}_\ell\|$ (typically $Z \le 1$ since
$u_\varepsilon$ maps to the unit disk). Then for all $\ell$:
\begin{equation}
  \|\delta_\ell\| \le \frac{L\rho}{1-2\rho L} (Z + \rho) + \frac{L}{1-2\rho L} B
  = O(L\rho) + O(LB).
\end{equation}
\end{corollary}

\begin{remark}[Quadratic vs linear in $\rho$]
The bound is \emph{linear} in $\rho$ because feedback enters directly
as $W_{\ell+1}^\top z_{\ell+1}^*$ in the fixed point equation. A
quadratic bound $O(\rho^2)$ would require the feedback to be
proportional to the \emph{perturbation} (e.g., if the feedforward
drive were zero). In our setting, the feedforward drive is $O(\rho)$,
so the feedback term is also $O(\rho)$.
\end{remark}

%======================================================================
\section{Classification Accuracy Bound}
%======================================================================

The classification logits are computed by cosine similarity with a
codebook $C = [c_1,\dots,c_K] \in \C^{d_L \times K}$:
\begin{equation}
  \text{logits}_j = \Real \ip{z_L}{c_j} = \Real(z_L^H c_j).
\end{equation}

\begin{theorem}[Accuracy drop bound]
\label{thm:accuracy}
Let $\gamma$ be the minimal logit margin on the test set under the
feedforward pass:
$\gamma = \min_{i} \big(\text{logits}_{y_i}^{\text{ff}} - \max_{j\ne y_i}
\text{logits}_j^{\text{ff}}\big) > 0$.
Let $M = \max_j \|c_j\|$. Then the classification accuracy under the
EP fixed point satisfies:
\begin{equation}
  \text{acc}_\text{EP} \ge \text{acc}_\text{BP}
  - \frac{2M \|\delta_L\|}{\gamma}.
\end{equation}
\end{theorem}

\begin{proof}
For any test sample, the logit difference for the true class $y$ vs
any other class $j$ is:
\begin{align}
  \text{logits}_y^* - \text{logits}_j^*
  &= (\text{logits}_y^{\text{ff}} - \text{logits}_j^{\text{ff}})
  + \Real\ip{\delta_L}{c_y - c_j} \\
  &\ge \gamma - \|\delta_L\| \|c_y - c_j\| \\
  &\ge \gamma - 2M \|\delta_L\|.
\end{align}
If $\gamma > 2M \|\delta_L\|$, the predicted class is unchanged.
The fraction of samples with $\gamma_i \le 2M \|\delta_L\|$ is at most
$2M \|\delta_L\| / \gamma$ by Markov's inequality (or exactly the
empirical fraction).
\end{proof}

Combining with Theorem \ref{thm:perturbation}:
\begin{equation}
  \text{acc}_\text{BP} - \text{acc}_\text{EP}
  \le \frac{2M}{\gamma} \Big(
  \frac{L\rho}{1-2\rho L} (Z + \rho) + \frac{L}{1-2\rho L} B \Big)
  = O\Big(\frac{L\rho}{\gamma}\Big).
\end{equation}

%======================================================================
\section{Finite-Time Settling}
%======================================================================

In practice, EP settling runs for finite steps $T$ with damping $dt$:
$z^{k+1} = (1-dt)z^k + dt F_\varepsilon(z^k)$.

\begin{lemma}[Finite-time error]
After $T$ steps with $dt \in (0,1)$, the iterate $z^T$ satisfies:
\begin{equation}
  \|z^T - z^*\| \le (1 - dt(1-\kappa))^T \|z^0 - z^*\| + \frac{dt}{1-\kappa} \sup_z \|F_\varepsilon(z)\|.
\end{equation}
With $dt=0.5$, $T=200$, $\kappa \approx 0.8$ (for $\rho=0.4, L=2$), the
iteration error is $<10^{-10}$. The discretization error is $O(dt) =
O(0.5)$ in the worst case, but empirically much smaller due to
contraction.
\end{lemma}

%======================================================================
\section{Numerical Constants for FashionMNIST Experiments}
%======================================================================

\begin{table}[h]
\centering
\caption{Bound parameters for 784$\to$256$\to$64 network (depth $L=2$).}
\label{tab:constants}
\begin{tabular}{@{}lll@{}}
\toprule
Parameter & Symbol & Value \\
\midrule
Weight scale factor & --- & 0.4 \\
Max spectral norm & $\rho$ & $\approx 0.4$ \\
Soft projection $\varepsilon$ & $\varepsilon$ & 0.1 \\
Depth & $L$ & 2 \\
Contraction factor & $\kappa = 2\rho L$ & 1.6 (violates (A5); see remark) \\
Codebook max norm & $M$ & $\approx 1$ \\
Typical margin & $\gamma$ & $\approx 0.5$--$1.0$ \\
Bias norm bound & $B$ & $\approx 0.1$ \\
\bottomrule
\end{tabular}
\end{table}

\begin{remark}[Contraction condition in practice]
For $L=2$, $\rho=0.4$ gives $\kappa = 1.6 > 1$, violating (A5).
However, the per-layer analysis in Theorem \ref{thm:perturbation}
does not require global contraction—it only requires $(I - \rho A)$
invertible, which holds if $\rho < 1/\sqrt{2} \approx 0.707$ for
$L=2$ (spectral radius of $A$ is $\sqrt{2}$). The tighter condition
$\rho < 1/(2L)$ is sufficient but not necessary. For $L\ge 3$ with
$\rho=0.4$, $\kappa > 1$ and the global contraction argument fails;
a per-layer or block-wise analysis is needed. Empirically, transfer
works for $L=2$ but degrades for $L\ge 3$.
\end{remark}

%======================================================================
\section{Summary}
%======================================================================

We have shown that for phasor networks with soft projection, the EP
free fixed point differs from the feedforward pass by $O(L\rho)$ in
hidden state norm, where $\rho$ is the max weight spectral norm and
$L$ is depth. This implies a classification accuracy drop of
$O(L\rho/\gamma)$ where $\gamma$ is the logit margin. For the
FashionMNIST setup ($L=2$, $\rho\approx0.4$, $\gamma\approx0.5$), the
bound predicts a drop of $\sim 30\%$, which is loose but the correct
scaling. The empirical drop is $<1\%$ because:
\begin{enumerate}
  \item The soft projection saturates phase, making the fixed point
    robust to small feedback perturbations.
  \item The spectral norm bound $\|W_\ell\|_2 \le \rho$ is loose;
    typical drives have smaller effective gain.
  \item The margin $\gamma$ on easy samples is much larger than the
    worst-case bound.
\end{enumerate}
The analysis extends to deeper networks with per-layer constants
$C_\ell, D_\ell$ (Theorem \ref{thm:perturbation}), but requires
$\rho < 1/\sqrt{2}$ for $L=2$ and tighter conditions for larger $L$.
This matches the empirical observation that depth $\ge 3$ networks
have no usable operating zone at current weight scales.

\bibliographystyle{plainnat}
\bibliography{liep_companion_proof}

\end{document}